% =========================================================
% Partiality, Totality, and Consistency
% =========================================================

\section{Maps and Order-Theoretic Notions}

We record several basic notions concerning partiality, totality, and compatibility, formulated in a way suitable for ordered structures.

% ---------------------------------------------------------
% Partial and Total Maps
% ---------------------------------------------------------

\begin{definition}[Partial and Total Maps]
Let $A$ and $B$ be sets.
\begin{itemize}
    \item A \textbf{partial map} $f : A \rightharpoonup B$ is a relation $f \subseteq A \times B$ such that for all $x \in A$ and $y_1,y_2 \in B$,
    \[
    (x,y_1) \in f \;\land\; (x,y_2) \in f \;\Longrightarrow\; y_1 = y_2.
    \]
    \item A \textbf{total map} $f : A \to B$ is a partial map such that for every $x \in A$ there exists $y \in B$ with $(x,y) \in f$.
\end{itemize}
Thus a total map is precisely a function in the usual sense.
\end{definition}

% ---------------------------------------------------------
% Consistency (Compatibility)
% ---------------------------------------------------------

\begin{definition}[Consistency / Compatibility]
Let $(S,\le)$ be a partially ordered set. Elements $x,y \in S$ are said to be \textbf{consistent} (or \textbf{compatible}) if they admit a common upper bound; that is,
\[
\exists z \in S \;\bigl(x \le z \;\land\; y \le z\bigr).
\]
\end{definition}

\begin{remark}[Order-theoretic interpretation]
Consistency expresses that $x$ and $y$ can be extended to a common refinement. In lattice-theoretic settings, this condition is automatic if binary joins exist, since one may take $z = x \vee y$.
\end{remark}

% ---------------------------------------------------------
% Maximal (Total) Elements
% ---------------------------------------------------------

\begin{definition}[Maximal Element]
Let $(S,\le)$ be a partially ordered set. An element $x \in S$ is \textbf{maximal} if
\[
\forall y \in S, \quad x \le y \;\Longrightarrow\; x = y.
\]
\end{definition}

\begin{remark}[Terminology]
In some contexts (notably domain theory), maximal elements are referred to as \emph{total objects}, reflecting the idea that they represent completed or fully specified states.
\end{remark}































% ---------------------------------------------------------
% Definition — Covering Relation
% ---------------------------------------------------------

\begin{definition}[Covering Relation]
Let $(P,\le)$ be a partially ordered set, and let $x,y \in P$ with $x < y$. We say that $y$ \textbf{covers} $x$ (or that $x$ is \textbf{covered by} $y$) if
\[
\forall z \in P, \quad (x \le z \le y \;\Longrightarrow\; z = x \;\text{or}\; z = y).
\]
Equivalently, $y$ covers $x$ if there is no element $z \in P$ such that
\[
x < z < y.
\]
\end{definition}

\begin{remark}[Finite posets]
If $P$ is finite, then for any $x,y \in P$ with $x < y$, there exists a finite sequence
\[
x = x_0 \prec x_1 \prec \cdots \prec x_n = y,
\]
where each $x_{i+1}$ covers $x_i$. Consequently, in a finite poset the order relation $\le$ is completely determined by the covering relation, and conversely the covering relation determines the order.
\end{remark}

% ---------------------------------------------------------
% Lemma — Characterizations of Order-Isomorphisms (Finite Case)
% ---------------------------------------------------------

\begin{lemma}[Characterizations of Order-Isomorphisms for Finite Posets]
Let $(P,\le_P)$ and $(Q,\le_Q)$ be finite partially ordered sets, and let $\varphi : P \to Q$ be a bijection. The following are equivalent:
\begin{enumerate}
    \item $\varphi$ is an order-isomorphism;
    
    \item for all $x,y \in P$,
    \[
    x <_P y \;\Longleftrightarrow\; \varphi(x) <_Q \varphi(y);
    \]
    
    \item for all $x,y \in P$,
    \[
    x \prec y \;\Longleftrightarrow\; \varphi(x) \prec \varphi(y),
    \]
    where $\prec$ denotes the covering relation.
\end{enumerate}
\end{lemma}

% ---------------------------------------------------------
% Proposition — Hasse Diagram Characterization
% ---------------------------------------------------------

\begin{proposition}[Hasse Diagram Characterization for Finite Posets]
Let $(P,\le_P)$ and $(Q,\le_Q)$ be finite partially ordered sets. Then $P$ and $Q$ are order-isomorphic if and only if their Hasse diagrams are identical up to relabeling of elements.
\end{proposition}

\begin{remark}[Interpretation]
In a finite poset, the order relation is completely determined by the covering relation. The Hasse diagram encodes precisely this covering structure. Consequently, two finite posets are order-isomorphic exactly when there exists a bijection between their elements that preserves the covering relation, equivalently, when their diagrams coincide up to relabeling.
\end{remark}

% ---------------------------------------------------------
% Definition — Dual of a Poset
% ---------------------------------------------------------

\begin{definition}[Dual of an Ordered Set]
Let $(P,\le)$ be a partially ordered set. The \textbf{dual} (or \textbf{opposite}) poset, denoted $P^{\mathrm{op}}$, is the set $P$ equipped with the order relation defined by
\[
x \le_{P^{\mathrm{op}}} y \;\Longleftrightarrow\; y \le_P x.
\]
\end{definition}

\begin{remark}[Diagrammatic interpretation]
If $P$ is finite, the Hasse diagram of $P^{\mathrm{op}}$ is obtained from that of $P$ by reversing all order relations. Equivalently, one may obtain it by turning the diagram of $P$ upside down.
\end{remark}

% ---------------------------------------------------------
% Principle — Duality
% ---------------------------------------------------------

\begin{proposition}[Duality Principle]
Let $\mathcal{S}$ be a statement about partially ordered sets expressed in terms of the order relation $\le$. If $\mathcal{S}$ holds in every partially ordered set, then the \emph{dual statement}, obtained by replacing each occurrence of $\le$ with $\ge$ (equivalently, reversing the order), also holds in every partially ordered set.
\end{proposition}

\begin{remark}[Interpretation]
The dual statement of $\mathcal{S}$ is obtained by systematically reversing all order-theoretic notions, such as exchanging upper bounds with lower bounds, suprema with infima, and maxima with minima. The principle reflects the fact that every partially ordered set $P$ has a dual $P^{\mathrm{op}}$, and any statement valid for all posets remains valid when applied to their duals.
\end{remark}

% ---------------------------------------------------------
% Definitions — Bottom and Top Elements
% ---------------------------------------------------------

\begin{definition}[Bottom and Top Elements]
Let $(P,\le)$ be a partially ordered set.

\begin{itemize}
    \item A \textbf{bottom element} (or \textbf{least element}) is an element $\bot \in P$ such that
    \[
    \forall x \in P, \quad \bot \le x.
    \]

    \item A \textbf{top element} (or \textbf{greatest element}) is an element $\top \in P$ such that
    \[
    \forall x \in P, \quad x \le \top.
    \]
\end{itemize}
\end{definition}

\begin{remark}[Uniqueness]
If a bottom element exists, it is unique. Likewise, if a top element exists, it is unique.
\end{remark}

\begin{example}[Bottom and Top Elements]
\begin{enumerate}
    \item In the poset $(\mathcal{P}(X), \subseteq)$, the bottom element is $\varnothing$ and the top element is $X$.

    \item Every finite chain has both a bottom element and a top element.

    \item An infinite chain need not have both. For example:
    \begin{itemize}
        \item The chain $(\mathbb{N}, \le)$ has a bottom element (namely $1$ or $0$, depending on convention), but no top element.
        \item The chain $(\mathbb{Z}, \le)$ has neither a bottom element nor a top element.
    \end{itemize}

    \item An antichain with more than one element has neither a bottom element nor a top element.
\end{enumerate}
\end{example}


\begin{remark}[Lifting]
Many results in order theory require the existence of a bottom element. Given a poset $(P,\le)$, one can adjoin a new element $\bot \notin P$ to obtain the \textbf{lifted poset} $P_\bot := P \cup \{\bot\}$, where the order is defined by
\[
\bot \le x \quad \text{for all } x \in P_\bot,
\]
and the original order on $P$ is preserved.

In particular, any set $S$ can be viewed as an antichain and then lifted to form a poset with a bottom element. Posets obtained in this way are called \textbf{flat}, since all elements of $S$ are incomparable above $\bot$.
\end{remark}


% ---------------------------------------------------------
% Definition — Lifting
% ---------------------------------------------------------

\begin{definition}[Lifting]
Let $(P,\le)$ be a partially ordered set. The \textbf{lifting} of $P$ is the poset $P_{\bot} := P \cup \{\bot\}$, where $\bot \notin P$, equipped with the order relation defined by
\[
\bot \le x \quad \text{for all } x \in P_{\bot},
\]
and
\[
x \le_{P_{\bot}} y \;\Longleftrightarrow\; x \le_P y \quad \text{for all } x,y \in P.
\]
Thus $P_{\bot}$ is obtained by adjoining a new bottom element to $P$.
\end{definition}

\begin{remark}
The element $\bot$ is the unique bottom element of $P_{\bot}$. This construction is often used to ensure the existence of a least element, for example in domain theory where $\bot$ represents an undefined or initial state.
\end{remark}

% ---------------------------------------------------------
% Definition — Co-Lifting
% ---------------------------------------------------------

\begin{definition}[Co-Lifting]
Let $(P,\le)$ be a partially ordered set. The \textbf{co-lifting} of $P$ is the poset
\[
P^{\top} := P \cup \{\top\},
\]
where $\top \notin P$, equipped with the order relation defined by
\[
x \le \top \quad \text{for all } x \in P^{\top},
\]
and
\[
x \le_{P^{\top}} y
\;\Longleftrightarrow\;
x \le_P y
\quad \text{for all } x,y \in P.
\]
Thus $P^{\top}$ is obtained by adjoining a new top element to $P$.
\end{definition}

\begin{remark}
The element $\top$ is the unique top element of $P^{\top}$. This construction is dual to lifting, which adjoins a bottom element.
\end{remark}


% ---------------------------------------------------------
% Definitions — Maximal and Minimal Elements
% ---------------------------------------------------------

\begin{definition}[Maximal and Minimal Elements]
Let $(P,\le)$ be a partially ordered set and let $Q \subseteq P$.

\begin{itemize}
    \item An element $a \in Q$ is \textbf{maximal in $Q$} if
    \[
    \forall x \in Q, \quad a \le x \;\Longrightarrow\; a = x.
    \]
    The set of maximal elements of $Q$ is denoted $\operatorname{Max}(Q)$.

    \item An element $a \in Q$ is \textbf{minimal in $Q$} if
    \[
    \forall x \in Q, \quad x \le a \;\Longrightarrow\; x = a.
    \]
\end{itemize}
\end{definition}

\begin{definition}[Greatest and Least Elements]
Let $(P,\le)$ be a partially ordered set and let $Q \subseteq P$.

\begin{itemize}
    \item An element $m \in Q$ is the \textbf{greatest element} (or \textbf{maximum}) of $Q$ if
    \[
    \forall x \in Q, \quad x \le m.
    \]
    If it exists, it is unique and denoted $\max Q$.

    \item An element $m \in Q$ is the \textbf{least element} (or \textbf{minimum}) of $Q$ if
    \[
    \forall x \in Q, \quad m \le x.
    \]
    If it exists, it is unique and denoted $\min Q$.
\end{itemize}
\end{definition}

\begin{remark}[Maximal vs.\ Maximum]
A maximal element need not be comparable with all elements of $Q$, whereas a greatest element must dominate all elements of $Q$. Thus:
\[
\max Q \;\text{exists} \;\Longrightarrow\; \operatorname{Max}(Q) = \{\max Q\},
\]
but the converse need not hold.
\end{remark}

\begin{proposition}[Finite Posets]
Let $(P,\le)$ be a finite partially ordered set.
\begin{itemize}
    \item Every nonempty subset $Q \subseteq P$ has at least one maximal element.
    \item For every $x \in P$, there exists $y \in \operatorname{Max}(P)$ such that $x \le y$.
\end{itemize}
\end{proposition}

\begin{remark}[Examples and Behavior]
\begin{itemize}
    \item A subset of $(\mathbb{N},\le)$ has a maximal element if and only if it is nonempty and finite.
    \item A subset may have many maximal elements, one, or none.
    \item For example, in $\mathcal{P}(\mathbb{N}) \setminus \{\mathbb{N}\}$, there is no greatest element, but each set of the form $\mathbb{N} \setminus \{n\}$ is maximal.
    \item The collection of all finite subsets of $\mathbb{N}$ has no maximal elements.
\end{itemize}
\end{remark}

\begin{remark}[Interpretation]
When the order models information, maximal elements often correspond to fully specified (or total) objects, while non-maximal elements represent partial or incomplete information.
\end{remark}

% =========================================================
% Maps Between Ordered Sets
% =========================================================

\section{Maps Between Ordered Sets}

This section introduces the structure-preserving maps between ordered sets. In particular, it provides the language needed to decide when two ordered sets have the same order-theoretic structure.

% ---------------------------------------------------------
% Definitions — Order-Preserving Maps, Embeddings, Isomorphisms
% ---------------------------------------------------------

\begin{definition}[Order-Preserving Map, Order-Embedding, and Order-Isomorphism]
Let $(P,\le_P)$ and $(Q,\le_Q)$ be partially ordered sets, and let
\[
\varphi : P \to Q
\]
be a map.

\begin{enumerate}
    \item The map $\varphi$ is \textbf{order-preserving} (or \textbf{monotone}) if, for all $x,y \in P$,
    \[
    x \le_P y \;\Longrightarrow\; \varphi(x) \le_Q \varphi(y).
    \]

    \item The map $\varphi$ is an \textbf{order-embedding} if, for all $x,y \in P$,
    \[
    x \le_P y \;\Longleftrightarrow\; \varphi(x) \le_Q \varphi(y).
    \]

    \item The map $\varphi$ is an \textbf{order-isomorphism} if it is an order-embedding and is onto $Q$.
\end{enumerate}
\end{definition}

\begin{remark}[Notation]
When $\varphi : P \to Q$ is an order-embedding, we write
\[
\varphi : P \hookrightarrow Q.
\]
When there exists an order-isomorphism from $P$ to $Q$, we say that $P$ and $Q$ are \textbf{order-isomorphic} and write
\[
P \cong Q.
\]
\end{remark}

% ---------------------------------------------------------
% Remarks — Basic Properties of Order-Preserving Maps
% ---------------------------------------------------------

\begin{remark}[Composition of Order-Preserving Maps]
Let $\varphi : P \to Q$ and $\psi : Q \to R$ be order-preserving maps. Then the composite map
\[
\psi \circ \varphi : P \to R,
\qquad
(\psi \circ \varphi)(x) = \psi(\varphi(x)),
\]
is order-preserving.

More generally, any finite composite of order-preserving maps is order-preserving, whenever the composite is defined.
\end{remark}

\begin{remark}[Image of an Order-Embedding]
Let $\varphi : P \to Q$ be an order-embedding, and let
\[
\varphi(P) := \{\varphi(x) : x \in P\}
\]
be its image. Then
\[
P \cong \varphi(P),
\]
where $\varphi(P)$ is ordered by the order inherited from $Q$.
\end{remark}

\begin{remark}[Order-Embeddings are Injective]
Every order-embedding is injective.

Indeed, if $\varphi : P \to Q$ is an order-embedding and $\varphi(x)=\varphi(y)$, then
\[
\varphi(x) \le_Q \varphi(y)
\qquad\text{and}\qquad
\varphi(y) \le_Q \varphi(x).
\]
By order reflection,
\[
x \le_P y
\qquad\text{and}\qquad
y \le_P x.
\]
By antisymmetry in $P$, it follows that $x=y$.
\end{remark}

\begin{remark}[Order-Isomorphisms and Inverses]
Ordered sets $P$ and $Q$ are order-isomorphic if and only if there exist order-preserving maps
\[
\varphi : P \to Q
\qquad\text{and}\qquad
\psi : Q \to P
\]
such that
\[
\psi \circ \varphi = \operatorname{id}_P
\qquad\text{and}\qquad
\varphi \circ \psi = \operatorname{id}_Q.
\]
\end{remark}

% ---------------------------------------------------------
% Lemma — Finite Characterizations of Order-Isomorphisms
% ---------------------------------------------------------

\begin{lemma}[Finite Characterizations of Order-Isomorphisms]
Let $(P,\le_P)$ and $(Q,\le_Q)$ be finite partially ordered sets, and let
\[
\varphi : P \to Q
\]
be a bijection. The following are equivalent:
\begin{enumerate}
    \item $\varphi$ is an order-isomorphism.

    \item For all $x,y \in P$,
    \[
    x <_P y
    \;\Longleftrightarrow\;
    \varphi(x) <_Q \varphi(y).
    \]

    \item For all $x,y \in P$,
    \[
    x \prec_P y
    \;\Longleftrightarrow\;
    \varphi(x) \prec_Q \varphi(y),
    \]
    where $\prec$ denotes the covering relation.
\end{enumerate}
\end{lemma}

\begin{proof}
The equivalence of (i) and (ii) follows immediately from the definitions.

Assume (ii), and suppose $x \prec_P y$. Then $x <_P y$, so by (ii),
\[
\varphi(x) <_Q \varphi(y).
\]
Suppose there exists $w \in Q$ such that
\[
\varphi(x) <_Q w <_Q \varphi(y).
\]
Since $\varphi$ is onto, there exists $u \in P$ such that $w=\varphi(u)$. By (ii),
\[
x <_P u <_P y,
\]
contradicting $x \prec_P y$. Hence
\[
\varphi(x) \prec_Q \varphi(y).
\]
The reverse implication is proved in the same way, using bijectivity of $\varphi$. Thus (iii) holds.

Now assume (iii), and suppose $x <_P y$. Since $P$ is finite, there exists a finite chain of covering relations
\[
x=x_0 \prec_P x_1 \prec_P \cdots \prec_P x_n=y.
\]
By (iii),
\[
\varphi(x_0) \prec_Q \varphi(x_1) \prec_Q \cdots \prec_Q \varphi(x_n).
\]
Therefore
\[
\varphi(x) <_Q \varphi(y).
\]
The reverse implication is proved similarly, using that $\varphi$ is onto. Hence (ii) holds.
\end{proof}

% ---------------------------------------------------------
% Proposition — Hasse Diagram Characterization
% ---------------------------------------------------------

\begin{proposition}[Hasse Diagram Characterization for Finite Posets]
Let $(P,\le_P)$ and $(Q,\le_Q)$ be finite partially ordered sets. Then $P$ and $Q$ are order-isomorphic if and only if their Hasse diagrams are identical up to relabeling of elements.
\end{proposition}

\begin{remark}[Interpretation]
For a finite poset, the order relation is determined by the covering relation. The Hasse diagram records exactly this covering relation. Thus two finite ordered sets are order-isomorphic precisely when their diagrams have the same covering structure, up to relabeling.
\end{remark}

% =========================================================
% The Duality Principle
% =========================================================

% ---------------------------------------------------------
% Definition — Dual Ordered Set
% ---------------------------------------------------------

\begin{definition}[Dual Ordered Set]
Let $(P,\le_P)$ be a partially ordered set. The \textbf{dual} ordered set, denoted $P^{\partial}$ or $P^{\mathrm{op}}$, has the same underlying set as $P$, with order relation defined by
\[
x \le_{P^{\partial}} y
\;\Longleftrightarrow\;
y \le_P x.
\]
\end{definition}

\begin{remark}[Diagrammatic Interpretation]
If $P$ is finite, a Hasse diagram for $P^{\partial}$ is obtained from a Hasse diagram for $P$ by turning the diagram upside down.
\end{remark}

% ---------------------------------------------------------
% Principle — Duality
% ---------------------------------------------------------

\begin{proposition}[Duality Principle]
Let $\Phi$ be a statement about ordered sets that is true for all ordered sets. Then the dual statement $\Phi^{\partial}$, obtained by reversing the order, is also true for all ordered sets.
\end{proposition}

\begin{proof}
Let $P$ be an ordered set. Since $P^{\partial}$ is also an ordered set, the statement $\Phi$ holds in $P^{\partial}$. Translating this statement back into the order of $P$ gives precisely the dual statement $\Phi^{\partial}$ in $P$.
\end{proof}

% =========================================================
% Down-Sets and Up-Sets
% =========================================================

\section{Down-Sets and Up-Sets}

Associated with any ordered set are two important families of subsets: down-sets and up-sets.

% ---------------------------------------------------------
% Definitions — Down-Sets and Up-Sets
% ---------------------------------------------------------

\begin{definition}[Down-Set and Up-Set]
Let $(P,\le)$ be a partially ordered set, and let $Q \subseteq P$.

\begin{enumerate}
    \item The subset $Q$ is a \textbf{down-set} if, whenever $x \in Q$, $y \in P$, and $y \le x$, one has
    \[
    y \in Q.
    \]
    Equivalently,
    \[
    x \in Q \;\land\; y \le x
    \;\Longrightarrow\;
    y \in Q.
    \]
    A down-set is also called a \textbf{decreasing set} or an \textbf{order ideal}.

    \item Dually, the subset $Q$ is an \textbf{up-set} if, whenever $x \in Q$, $y \in P$, and $x \le y$, one has
    \[
    y \in Q.
    \]
    Equivalently,
    \[
    x \in Q \;\land\; x \le y
    \;\Longrightarrow\;
    y \in Q.
    \]
    An up-set is also called an \textbf{increasing set} or an \textbf{order filter}.
\end{enumerate}
\end{definition}
